\section{Readouts and pre-registered predictions}\label{sec:readouts}

The kernel of Section~\ref{sec:kernel} has no dynamics, so numbers can only enter as \emph{readouts}: identifications of a structural
quantity with a measured one.  This section lists the readouts that survive the audits of the Companion, and the four predictions
that were fixed before the relevant data appeared.  Since October 2026 the project has followed one rule: every test is written down with
its data and decision criterion and committed \emph{before} the computation; notes that were not blind say so (Appendix~\ref{app:method}).

\subsection{The matter fraction}

The Companion identifies the density fractions with the frame offset:
\[
\Omega_M=\beta_{\rm sym}^2=\frac{\gamma_M-\gamma_T}{\gamma_M+\gamma_T}=0.3018,\qquad \Omega_\Lambda=1-\beta_{\rm sym}^2=0.6982,
\]
where the second form uses the two countings of Proposition~\ref{prop:countings}, \(\gamma_T=\gamma_{\rm string}\) and \(\gamma_M=0.2375330\)
(layer v4.9.0-r67).  For comparison, DESI DR2 BAO alone gives \(\Omega_m=0.2975\pm0.0086\)~\cite{desi} (\(+0.5\sigma\)) and Planck~2018
\(0.3153\pm0.0073\)~\cite{planck} (\(-1.8\sigma\)).  Three restrictions apply.
\begin{itemize}
\item \(\beta_{\rm sym}\) is constant while \(\Omega_M\) changes with cosmic time.  The identification can hold at one epoch only; at all epochs it
would give \(\Omega_\Lambda=0.698\) at recombination, which is excluded.  It is a statement about the present epoch, and why now is open.
\item What the identification concerns is the \emph{share} \(\Omega_\Lambda\), not \(\Lambda\) itself.  Under the counting reading the frames differ
in their counting, not in \(G_N\) or the geometry, so \(\Lambda\) as a curvature is frame-independent.  The older claim that a matter observer
measures a cosmological constant \emph{because of} his frame offset is not supported, and two remarks of the Companion that put
\(\beta_{\rm sym}^2\) on \(\Lambda\) instead of on \(\Omega_M\) contradict its own theorem.
\item The comparison is after the fact; it was never registered as a prediction.
\end{itemize}

\subsection{Mass forms}

The charged-lepton masses satisfy the Koide relation~\cite{koide}
\[
K=\frac{m_e+m_\mu+m_\tau}{\bigl(\sqrt{m_e}+\sqrt{m_\mu}+\sqrt{m_\tau}\bigr)^2}=\frac23
\]
to about \(10^{-5}\).  This is an empirical fact; its derivation from the triple structure is open.  What the framework does give is the
\emph{form}.  In the operator language of Section~\ref{sec:generation}, mass terms couple the matter frame to the tachyon frame and carry
the quantum numbers of a Higgs doublet (layer v4.9.0-r44).  A coupling with the full slot symmetry (cyclic shift \(\tau\) and reflection
\(\rho\)) gives at most two distinct masses; breaking \(\rho\) and keeping \(\tau\) gives a circulant coupling \(a+bP+\bar bP^2\) whose
eigenvalues are three distinct masses of the form
\[
\sqrt{m_k}=a+2|b|\cos\Bigl(\theta+\frac{2\pi k}{3}\Bigr),\qquad k=0,1,2,
\]
and the Koide value corresponds to \(|b|=a/\sqrt2\).  Mixing needs the breaking of \(\tau\) itself.  So each observed property has its breaking
stage --- mass the frame coupling, mass splitting the breaking of \(\rho\), mixing the breaking of \(\tau\) --- but neither \(|b|/a\) nor
\(\theta\) is explained, and many direct attempts to derive the masses from natural potentials were refuted (pre-registered tests H11,
H16--H18).

\subsection{A mixing model}

For the lepton mixing angles the Companion uses a readout rule: each particle kind reads its mass coupling with a subgroup of
\(A_4\supset\mathbb Z_3=\langle\tau\rangle\).  Dirac couplings are circulant and read with \(\mathbb Z_3\); Majorana neutrinos admit a further
coupling and read with a \(\mathbb Z_2\).  This gives the known trimaximal pattern TM2 with \(\sin^2\theta_{12}=1/(3\cos^2\theta_{13})\approx0.341\),
which lies \(3.6\sigma\) above JUNO~\cite{juno}.  A correction \(D=\exp(\varepsilon W)\) with a fixed diagonal \(W\) (the ``mask'') lowers it into
the measured range while \(\theta_{13}\) and \(\theta_{23}\) stay inside theirs (layers v4.9.0-r55, r57).  Since \(D\) is boost-like, the
natural value of \(\varepsilon\) is a rapidity, and the reading \(\varepsilon=\eta=\operatorname{artanh}\beta_{\rm sym}\approx0.617\) meets the
data, while the Hopf value \(\operatorname{artanh}(1/\sqrt3)\) does not (layer v4.9.0-r60).  This model is a fit with few settings, built on the
\(A_4\) literature; the reading \(\varepsilon=\eta\) was chosen after the target windows were known.  It is the source of predictions H42 and
H46 below.

\subsection{Pre-registered predictions}

\begin{center}
\small
\renewcommand{\arraystretch}{1.3}
\begin{tabular}{@{}>{\raggedright\arraybackslash}p{0.12\textwidth}>{\raggedright\arraybackslash}p{0.31\textwidth}>{\raggedright\arraybackslash}p{0.24\textwidth}>{\raggedright\arraybackslash}p{0.23\textwidth}@{}}
\toprule
& prediction & data that count & refuted if\\
\midrule
H9 (20 Oct 2026) & \(\sin^2\theta_{12}=\beta_{\rm sym}^2/\cos^2\theta_{13}=0.3086\pm0.0002\) & next JUNO result with \(\sigma\le0.004\) & pull \(|z|\ge3\)\\
H15 (21 Oct 2026) & \(\cot\theta_W=\gamma_L\) (on shell), i.e.\ \(m_W=m_Z\gamma_L/\sqrt{1+\gamma_L^2}=80.3615\pm0.0018\) GeV & first world average or official combination after 21 Oct 2026 & criterion in the note\\
H42 (24 Oct 2026) & \(\sin^2\theta_{23}\le0.54\) (normal ordering) & first global fit or DUNE/Hyper-K result after 24 Oct 2026 & \(2\sigma\) interval wholly above \(0.54\)\\
H46 (24 Oct 2026) & \(\sin^2\theta_{23}\in[0.49,0.50]\) under \(\varepsilon=\eta\) & as H42 & \(2\sigma\) interval misses \([0.49,0.50]\)\\
\bottomrule
\end{tabular}
\end{center}

Disclosures, as recorded in the notes.  H9 and H15 were noticed during other computations and lie close to known values (H9 \(0.7\sigma\)
from the combination of JUNO and NuFIT~6.0, H15 \(0.8\sigma\) from the on-shell value); data known at the time do not count, only the next
qualifying result.  H42 was noticed after an evaluation and holds as a prediction only from its note on.  H46 rests on a reading that was
not blind, although the value of \(\theta_{23}\) was not known when the reading was chosen.  Present global fits do not decide H42 or
H46: NuFIT gives \(\sin^2\theta_{23}=0.561\) without and \(0.470\) with Super-Kamiokande atmospheric data~\cite{nufit}.  A confirmation of
H46 would test the chain of the mixing model, not the reading alone, since other models also predict \(\theta_{23}\) near \(45^\circ\).
To our knowledge no qualifying data had appeared when this was written (3 November 2026).

Two earlier structural checks belong here.  Neutrinos as a light-like triple give normal ordering and \(\sum m_\nu=0.0592\) eV (H2); normal
ordering is supported by JUNO.  And of the six predictions P1--P6 stated in the original Companion, none is confirmed: P1 is falsified, P2
excluded, P3 a coincidence, P4 contradictory, P5 corrected, P6 untested (Section~\ref{sec:not}).

\medskip\noindent\emph{Status.}  All readouts in this section are identifications, not derivations; they inherit the input status of
\(\beta_{\rm sym}\).  The four predictions are falsifiable and pending.  Their value lies less in their chance of being right than in being
fixed in advance.
